\documentclass[dvipdfmx]{jlreq}

\title{「粘土による形状伝達,多義プロンプトの解釈提案と制約探索を用いたデザイン支援手法」章立て}
\author{梅木敦成}
\date{\today}

\begin{document}
\maketitle


% タイトルに入れるぐらい重要な用語なら「多義プロンプト」はどこかでちゃんと説明しないとダメですね．１章でも登場するだろうし，アプローチの概要でもう一回ちゃんと話すのか．ok
% 「解きたい課題」「なぜ現状の他の手法ではうまく行っていないのか」が不足している
% 「なら生成ai」の前に，一般的にどういう支援がこれまで考えられてきたのか，それに対して「生成ai」はどういう良さがあると考えるのか
% その上で，「生成ai」をデザイン生成に利用している多くの人は実際どうやって解決しているのか，どういう形で妥協，というか成果物を得ているのか，それに対する問題はなんなのか
% 「出ない」止まりはまだ弱いかなok
% システムの提案がユーザの想定外という話は突然出てきたがどこから？ok
% 先行研究の話を盛り込むok


\section{はじめに}
\subsection{研究背景}
\begin{itemize}
  \item ものづくり、デザインでの創作が背景
  \item 非専門家にとって頭の中のイメージをカタチにするのは簡単ではない
  \item 従来はスケッチやCADツールが用いられてきたが、これらは高い操作スキルや描画能力を必要とするため、非専門家には障壁が高かった。
  \item しかし、近年 Text-to-Image 生成AIが登場したことで、この「描画スキルの壁」は取り払われた。 言語（プロンプト）という誰もが使えるインターフェースにより、造形への参入障壁は劇的に下がった。
  % \item なら生成aiを用いた情報システムで支援してやろう
  % \item 近年、Text-to-Image生成AIが登場し、技術的な障壁は下がった(言語（プロンプト）という、誰もが使えるインターフェースで操作できる点に革新性があった)。
  \item 多くのユーザーはプロンプトエンジニアリング（試行錯誤）によって解決を図っている。
  \item しかし、非専門家は専門的な造形言語（形状や質感の名称）を持たないため、「かっこいい」「強そうな」といった抽象的で多義的な語彙（以下、多義プロンプトと呼ぶ）を用いて指示を行う傾向がある。
  \item 多義プロンプトは解釈の幅が広いため、AIが学習データに基づき出力する結果と、ユーザー個人の文脈やこだわり（メンタルモデル）との間に「解釈の乖離」が生じやすい。
  \item 単に生成を繰り返しても、ユーザーは「どの言葉を変えれば意図が伝わるのか」が分からず、偶然の正解が出るまで試行錯誤を強いられる（＝意図通りのデザインが出力されない、効率が悪い）。
  \item また、「かっこいい」などの雰囲気（スタイル）は言葉で伝えやすいが、「ここが少し膨らんでいて、あそこが凹んでいる」といった具体的な「形状」の部分修正をテキストだけで正確に伝えることは不可能に近い。(部分修正/こだわりの表現が難しい)
  \item 既存の Image-and-Text-to-Image でオンラインの画像で検索して補完する手法もあるが、ユーザーがイメージする形状特徴を持つ画像が必ずしも存在するとは限らない。
  \item 先行研究であるStyleCLIPなどは、テキスト入力によって画像の特徴を操作する手法を提案している。しかし、これらは計算機的な整合性（テキストといかに一致するか）を重視しており、ユーザー個人の『メンタルモデルとの整合性』や『主観的な満足度』を直接支援するものではない。そのため、非専門家が自身の曖昧なイメージを具体化する用途には適していない。
  \item そこで本研究では、生成AIにおける「言語化の壁」を突破するために、言語化しにくい形状特徴を直接入力できる（粘土などの）物理的なインターフェースと、AIによる解釈支援を組み合わせた手法を提案する。
  % \item しかし非専門家が単にプロンプティング頑張って生成aiを使ってもユーザーにとって「メンタルモデルとの整合性」（ふわふわした語彙で良ければ「納得感,満足感」)の高いデザインはなかなか出ない
  % \item 単に生成aiを使ってもユーザーの意図通りのデザインはなかなか出ない理由の一つは専門的な知識を持たないユーザーは抽象的な、多義的な言葉でプロンプティングしてしまう傾向があり、aiがプロンプトをユーザーの意図とは異なる解釈をしてしまうこと。
  % \item 本研究はユーザーの曖昧的なイメージを、システムとの対話を通じて段階的に具体化し、要求仕様を洗練する支援システムを提案する。
  % \item ユーザーがイメージしたデザインと生合成の高いデザインの生成を支援することを目指す。
  % \item また、システムのする提案がユーザーの想定外のであることでユーザーの発想が広がることも目指す
\end{itemize}

\subsection{本研究のアプローチと目的}
\begin{itemize}
  \item 本研究では、ユーザーの曖昧なイメージを、システムとの対話を通じて段階的に具体化し、要求仕様を洗練する支援システムを提案する。
  \item 具体的には、多義プロンプトを一義的な解釈候補へと具体化し、ユーザーの潜在的な意図を言語化・明確化し、粘土による形状伝達と組み合わせることで、ユーザーの意図に沿ったデザイン生成を支援する。
  \item この探索プロセスを実現するために、本研究では検索支援システムFindMeで提案された『制約に基づくナビゲーション』の概念を応用する。FindMeは既存のカタログからの検索を対象としていたが、本システムではこれを生成モデルの潜在空間探索へと拡張し、『もっと〇〇に』という比較制約による直感的なデザイン生成を実現する。
  \item また、意図した範囲内（制約内）で多様な候補を提示することで、ユーザー自身も気づいていなかった好みの発見（発想の拡張）を促す。
  \item 最終的に、非専門家ユーザーのメンタルモデルとの整合性が高いデザイン生成を支援することを目的とする。

  % \item しかし非専門家が単にプロンプティング頑張って生成aiを使ってもユーザーにとって「メンタルモデルとの整合性」（ふわふわした語彙で良ければ「納得感,満足感」)の高いデザインはなかなか出ない
  % \item 単に生成aiを使ってもユーザーの意図通りのデザインはなかなか出ない理由の一つは専門的な知識を持たないユーザーは抽象的な、多義的な言葉でプロンプティングしてしまう傾向があり、aiがプロンプトをユーザーの意図とは異なる解釈をしてしまうこと。
  % \item 本研究はユーザーの曖昧的なイメージを、システムとの対話を通じて段階的に具体化し、要求仕様を洗練する支援システムを提案する。
  % \item ユーザーがイメージしたデザインと生合成の高いデザインの生成を支援することを目指す。
  % \item また、システムのする提案がユーザーの想定外のであることでユーザーの発想が広がることも目指す
\end{itemize}


% 「合意を取る」というより，「ユーザの曖昧な興味・状態を明確化・言語化する」みたいな方向が重要なのでは
% この時点では「特徴ベクトル」みたいなとこまで言わなくても良いのでは．２章の方で具体化で良いのでは
% 「さらに制約を用いた潜在空間を狭めることで細部の微調整を行う。」←潜在空間とは？
% 「システムは狭まった潜在空間中のランダムサンプリングで画像を生成するので」←これは突然　どこから来た？

\subsection{提案手法の概要}
\begin{itemize}
  \item 粘土により大まかな特徴（形状）を伝達する。
  % \item システムは多義的なクエリを具体的な視覚的特徴（属性）の候補に分解して提示する。
  \item システムは多義的なクエリに対し、それを具体化した複数の「解釈（視覚的特徴のセット）」を候補として提示する。
  \item ユーザーはその中から自身のイメージに近い解釈を選択することで、曖昧だった意図をシステムに対して明確化（言語化）し、探索の初期地点（特徴ベクトル）を決定する。
  \item さらに、生成画像の変種したい部位に対し「もっと〇〇に」といった比較形式の制約を与えることで、探索範囲をユーザーの好みの領域へと絞り込む。
  \item システムは、この絞り込まれた領域（潜在空間：デザインの特徴が地図のように配置された空間）の中から、多様なバリエーションをランダムに提示する。
  \item これにより、ユーザーの意図（制約）を守りつつ、ユーザーが想像していなかった細部のデザイン（想定外の解）に出会う機会を提供し、発想を広げる。
  % \item クエリの解釈を属性に分解し、ユーザーに確認することでまず大まかの方向性について合意をとる。
  % \item ユーザが合意した解釈に基づき特徴ベクトルを生成
  % \item さらに制約を用いて潜在空間を狭めることで細部の微調整を行う。
  % \item ここで潜在空間とは、デザインの特徴を表すベクトルが配置された空間を指す。この空間では「意味的な特徴（丸い、赤い、金属）」が数値化されて座標となっている。だからこそ、この空間内で数値をいじると、画像の形を保ったまま『質感だけ変える』といった操作が可能になる。
  % \item システムは狭まった潜在空間中のランダムサンプリングで画像を生成するので、
  % \item ユーザーにとってメンタルモデルとの整合性が高く、かつユーザーの想定外のデザインが生成されることでユーザーの発想が広がるのではないか
\end{itemize}

% \subsection{関連研究}
% \begin{itemize}
%   \item findmeの制約によるによる特徴ベクトル空間の切断を利用
%   \item StyleCLIPは特徴ベクトルの属性値を操作して画像の細部を調整する点で似ているがその研究は、テキストプロンプトとの整合性を機械的に評価し、スコアの向上を目指すのに対し、本研究は人間であるユーザーがデザインを評価し、ユーザの「メンタルモデルとの整合性」（ふわふわした語彙で良ければ「納得感,満足感」）の向上を目指す。
% \end{itemize}









\section{提案手法}
\subsection{アプローチの概要}
\begin{itemize}
  \item 非専門家のユーザーは、自身のイメージを表現するための専門的な語彙を持たないため、どうしても「かっこいい」のような多義的で抽象的な語彙に頼らざるを得ない。(繰り返し)
  \item しかし、このような抽象的なクエリは、解釈の多様性からユーザの意図しない画像の生成に繋がりうる。
  \item そこで本研究では、この「抽象的な入力による解釈の乖離」を解消し、意図を具体化するために、以下の2段階のアプローチを提案する。
\end{itemize}

% \subsection{粘土による形状入力の拡張}
% \begin{itemize}
%   \item ユーザーが作成した粘土のラフモデルを画像生成AIへの形状の指示情報として利用する。
% \end{itemize}

\subsection{形状入力と多義的クエリの解釈提案による合意}
\begin{itemize}
  \item まずユーザーが作成した粘土のラフモデルを画像生成AIへの形状の指示情報として利用する。
  \item その後抽象的なクエリが持つ多義性を解消するために、「解釈の提案」を行う。
  \item 非専門家の入力するクエリ（例：「和風」）は、多様な解釈が可能である。
  \item システムは、そのクエリに関連する代表的な視覚的特徴（属性の組み合わせ）を複数パターン生成し、ユーザーに提示する。
  \item 例：解釈A：[属性：茶色、低彩度、木質]（落ち着いた和風）, 解釈B：[属性：赤、白、金]（祝祭的な和風）
  \item 属性：「金属的」「角が丸い」といった具体的な特徴
  \item 属性値：その適合度合いを表す0 から1 の実数値。
  \item ユーザーがいずれかの解釈を選択することで、システムは予め設計した6グループの属性空間を基にユーザーが求めている「和風」の具体的なパラメータ（特徴ベクトル）を特定し、方向性の合意形成を行う。
\end{itemize}

\subsection{制約に基づく探索}
\begin{itemize}
  \item 次に、「制約」を用いた探索によって、形状や細部のニュアンスを調整する。
  \item 方向性が定まった後も、細部の調整においては言語化が難しい微細な意図が存在する。
  \item 生成された画像を見ながら、ユーザーが特定の属性に対して「制約」を追加する。
  \item 制約：属性値P$\le$0.6のように不等号を用いて属性の適合度合いを制限する
  \item システムはユーザーが追加した制約を満たす空間内で特徴ベクトルを更新
\end{itemize}

\subsection{システム構成と動作}
\begin{itemize}
  \item 使用した技術：Streamlit, GPT-4o, GPT-Image-1, Inpainting API
  \item 6グループの属性空間を設計
  \item 最初にモードを選択
  \item 全体編集モード
  \begin{itemize}
    \item (A)初期入力：粘土画像＋粘土が表す物体名入力＋自由テキスト入力（多義的クエリ）→ (B)
    \item (B)解釈の提案：システムは(A)のクエリに対応する複数の解釈を提示し、ユーザーが選択する。この解釈を元に生成AIを用いて特徴ベクトルを生成する。→ (C)
    \item (C)画像生成：特徴ベクトルを用いて画像を生成し、ユーザーに提示する。→ (J)
    % \item (D)制約付与と生成：生成画像に対し、ユーザーは特定の属性について制約を追加する。システムは制約を満たす部分空間内で特徴ベクトルを更新し、画像を再生成する。ユーザーが満足するまで(C)を繰り返す。
  \end{itemize}
  \item 部分編集モード
  \begin{itemize}
    \item (D)初期入力：粘土画像 ＋ キャンバス描画によるマスク作成 ＋ 粘土が表す物体名入力 ＋ 編集パーツ名/部位名入力 ＋ 自由テキスト入力（多義的クエリ）→ (E)
    \item (E)解釈の提案：システムは(D)のクエリに対応する複数の解釈を提示し、ユーザーが選択する。この解釈を元に生成AIを用いて部分特徴ベクトルを生成する。→ (F)
    \item (F)画像生成：特徴ベクトルを用いて画像を生成し、ユーザーに提示する。→ (J)
  \end{itemize}
  \item 合成モード(パーツを本体に合成する)
  \begin{itemize}
    \item (G)初期入力：粘土2パーツを写した画像 ＋ キャンバス描画によるマスク作成 ＋ 粘土が表す本体名入力 ＋ 合成パーツ名入力 → (H)
    \item (H)解釈の提案：システムは(G)の合成要求をユーザーに確認。→ (I)
    \item (I)画像生成：合成した画像を生成し、ユーザーに提示する。→ (J)
  \end{itemize}
  \item (J)制約付与と生成：生成画像に対し、ユーザーは特定の属性について制約を追加する。システムは制約を満たす部分空間内で特徴ベクトルを更新し、画像を再生成する。ユーザーが満足するまで繰り返す。
  \begin{itemize}
    \item 全体調整：生成画像に対し、ユーザーは特定の属性について制約を追加する。システムは制約を満たす部分空間内で特徴ベクトルを更新し、画像を再生成する。
    \item 部分調整：(D)〜(F)
    \item 部分微調整：すでに部分ベクトルのあるパーツ/部位の調整にはベクトルの再生成は不要。制約付与→特徴ベクトル更新→画像再生成。
  \end{itemize}
  \item 保存して終了/初期入力に戻る
  \item ユーザーは終了時の画像を参考に手元の粘度模型を修正し、再度システムに入力することで、反復的にデザインを洗練できる。
  \item 補足：斥力による多様性維持
  \begin{itemize}
    \item 特徴ベクトル空間内で、既に生成されたデザインのうちユーザーが「悪い」と判断したデザインに近い領域を避けるように調整を行うことで、多様なデザインの生成を促進する。
  \end{itemize}
\end{itemize}



\section{実験例と考察}
本章では、提案手法（形状入力と解釈提案と制約探索）が、デザインの専門知識を持たないユーザーにとって、メンタルモデルとの整合性が高い画像を生成するうえで有効であるかを検証するための実験について述べる。具体的には、従来のプロンプト入力のみの手法と比較し、探索の効率性（定量的評価）および、ユーザーの主観的な満足度や発想支援効果（定性的評価）が向上することを確認することを目的とする。
\subsection{実験方法}
\begin{itemize}
  \item タスク設定：正解画像を用意せず抽象的なお題を3題用意する。実用性を考慮し、日用品を対象とする。
  \item 題1「かっこいい椅子をデザインしてください」
  \item 題2「かわいいカップをデザインしてください」
  \item 題3「」
  \item 各条件でユーザーが満足するまで生成、時間無制限
  \item 提案手法: 形状入力とテキストによる初期入力。解釈提案と特徴ベクトル生成を経て、制約による微調整を行いながら反復的に画像生成
  \item 比較対象設定1: 形状入力とテキストによる初期入力。テキストプロンプトのみで画像反復生成。反復して画像を生成する際直前の生成デザインを入力。
  \item 提案手法と比較し、「特徴ベクトル」や「制約」がないと微調整で全体が崩れてしまい、制御性が低いことを証明
  \item 比較対象設定2: テキストによる初期入力のみ。テキストプロンプトのみで画像反復生成。反復して画像を生成する際直前の生成デザインを入力しない。
  \item 比較対象1と比較し、て形状入力がないとユーザーの意図する形状特徴を反映しにくいことを証明
  \item 定量的評価：座標の推移,収束までに生成した画像枚数などを計測し、効率性を比較する
  \item 定性的評価：作業中や作業後のインタビュー、作業中の発言内容の分析により、システムの使用感、デザインの満足度を測る
\end{itemize}
\subsection{結果の説明}
\subsection{考察}





\section{結論}
\subsection{この論文では何を述べたか}
\begin{itemize}
  \item 形状入力と、生成AIを用いたクエリの解釈提案と制約を用いた探索空間の制限によって非専門的なユーザーでもメンタルモデルとの整合性の高いデザインの創作が支援され、簡単になることを目指した
\end{itemize}
\subsection{今後の課題}
結果考察から得られるシステムの改良点


\end{document}
